\documentclass[]{article}

\usepackage{amsmath}
\usepackage{multirow}
\usepackage{natbib}
\usepackage{graphicx} 


\begin{document}

% Auto-generated figures file
% \usepackage{graphicx}

\begin{figure}[h!]
    \centering
    \includegraphics[width=0.5\textwidth]{../input_files/plots/step_5_IR_Decay_Combined_1_20260406_002623.png}
    \caption{Performance comparison of portfolio construction methodologies for the low-Sharpe SMB (left panels) and high-Sharpe WML (right panels) factors as a function of cross-sectional size (N). The top panels plot the Annualized Information Ratio (IR) for Ideal Factor Mimicking Portfolios (Ideal FMP), OLS-estimated FMPs (OLS-FMP), and Characteristic-Sorted portfolios. For the low-Sharpe SMB factor, the OLS-FMP performance attenuates with increasing N due to idiosyncratic noise leakage, whereas the Characteristic-Sorted portfolio remains more stable. Conversely, for the high-Sharpe WML factor, the OLS-FMP successfully captures the premium for N $\ge$ 20, while the Characteristic-Sorted approach fails. The bottom panels display the IR difference between the OLS-FMP and Characteristic-Sorted portfolios ($\Delta$IR = IR\_{OLS} - IR\_{Char}) with 95\% bootstrap confidence intervals, showing that OLS-FMP is statistically superior for WML, while the two methods become statistically indistinguishable for SMB at N $\ge$ 40.}
    \label{fig:image0}
\end{figure}

\begin{figure}[h!]
    \centering
    \includegraphics[width=0.5\textwidth]{../input_files/plots/step_5_Noise_to_Signal_Ratio_2_20260406_002623.png}
    \caption{Evolution of the Noise-to-Signal Ratio (NSR) on a logarithmic scale for OLS-estimated Factor Mimicking Portfolios (OLS-FMP) as a function of cross-sectional size (N). For the low-Sharpe SMB factor, the NSR demonstrates a catastrophic increase from 263.8 at N=10 to 6572.7 at N=50, indicating that idiosyncratic volatility leakage from noisy beta estimates overwhelms the factor signal as the cross-section expands. In contrast, the high-Sharpe WML factor shows signal robustness, with its NSR improving from 511.7 to 301.9, highlighting that a strong true premium provides a buffer against estimation error degradation.}
    \label{fig:image1}
\end{figure}



\end{document}
                